\chapter{Project Objectives}\label{ch:obiective}
\pagestyle{fancy}


\section{Vision}

OpenCB aims to be a software solution, designed for chess players of all skill levels, which serves as an add-on for traditional chess, to assist when playing. Checking the validity of chess movements, OpenCV is a great tool for beginners who struggle to respect the rules of chess, but also helpful for professional players, since we all make mistakes. A key idea while developing the project is to keep the physical aspect of the game intact as much as possible. At the same time, the project is open source, under an MIT license, hence available for any developer willing to use the code or contribute to it.

\b{OpenCB is a mixture of hardware and software. In the hardware part, cameras are used to get images of a physical board. In the software part, Image Processing algorithms are applied to classify the pieces seen by the cameras and custom algorithms are used to validate the movement we get from the reconstructed boards.}



\section{Users and Use Cases}

From the Vision we can derive that we have two types of users. The main one will simply be called \b{User}. By this, we refer to people who download the software, set it up and want to use it. One such user can use the system for many use cases, which can be grouped into categories:

\begin{itemize}
    \item Chess piece image acquisition
    
    We need to get images of chess pieces. To achieve this, a previous configuration will be necessary, to get the corners of the board (\i{Configure}), after which these cell images can be extracted and labeled (\i{Crop and Label}). To be more user-friendly, test variations of these functions shall be provided as well, such that the results of a potential Configure or Crop action can be seen without any effects on previous values.
    
    \item Chess piece image classification

    After having the piece images, we can use Machine Learning algorithms to classify them as either White Free, Pawn, Bishop, Knight, Rook, Queen, King or one of the black counterparts. For this, multiple types of classifiers will be defined, from which the user can choose (\i{Select Classifier}). These can be saved or loaded to save time, trained if necessary and tested to check for metrics. The main functionality \i{Classify Board} can then take a set of 8x8 cell images and classify all of them, reconstructing a board.

    \item Move validation
    
    If we have a board position, we can compare it to the previous one, and see if a valid move happened ``in between''. This is done by \i{Validate Move}. Some moves might need to be discarded, as well as there needs to be an easy way of starting a new game.

    \item Miscellaneous actions
    
    Aside from what has already been listed, the User can also \i{Shuffle and Split} to mix up the cell images, \i{Clear All} to delete all images, or \i{Change Settings} to modify parameters which persist cross-session.
\end{itemize}

The secondary actor of the system is \b{Developer}. This could be any open-source connoisseur willing to contribute to the codebase or reuse parts of it. For this, it is essential to have readable code, good documentation and clear version control.

These two actors and the associated use cases can be seen on Figure \ref{fig:usecase}.


\begin{figure}[H]
    \centering
    \includegraphics[width=0.6\textwidth]{figs/usecase.png}
    \caption{Use Case Diagram}
    \label{fig:usecase}
\end{figure}



\section{Other Product Requirements}

As stated before, readability and documentation are key. Another important aspect is performance since it directly impacts the user experience. There are parts of the project that are real-time, like providing a preview of what the camera sees. Here we need an FPS high enough not to annoy the user. Other times, like when classifying, the user can wait a little, since he/she knows that processing is going on in the background, but the wait should not be longer than necessary.

Maintainability and extendibility are also important since the project might live on for longer. There are a lot of ideas not implemented up to this point, and a lot of other ideas might come in the future. By trying not to overcomplicate solutions, commenting where necessary and leaving space for future additions, the project has a greater value as code.
